\chapter{Etude préalable}

\section{Description de l'existant}
L'étude de l'existant a mis en évidence un besoin de structurer les interfaces par rôle utilisateur, avec des parcours clairs pour les opérations quotidiennes (cours, devoirs, notes, absences, demandes).

\section{Critiques de l'existant}
Les principales limites identifiées étaient :
\begin{itemize}
	\item navigation peu intuitive ;
	\item hétérogénéité visuelle entre modules ;
	\item manque de lisibilité sur les états métier ;
	\item interactions API non uniformes côté client.
\end{itemize}

\section{Solutions proposées}
La solution proposée consiste à développer une SPA React/TypeScript avec :
\begin{itemize}
	\item routage structuré et protection des accès ;
	\item composants réutilisables ;
	\item pages spécialisées par rôle ;
	\item appels API factorisés et robustes.
\end{itemize}

\section{Développement de la solution finale}
La solution finale frontend fournit des tableaux de bord fonctionnels et des écrans opérationnels pour les profils enseignant et étudiant, avec intégration des services backend et amélioration de l'ergonomie globale.
